\chapter*{Abstract}
\addcontentsline{toc}{chapter}{Abstract}

This report documents the design, implementation, and validation of the LEO Rover Mission Control system, a full-stack autonomous mission control platform developed during an engineering internship at the Florida Institute of Technology's robotics laboratory, under the supervision of Dr.\ Gutierrez, as part of the Carolus multi-robot indoor-localization research program.

At the outset of the internship, the inherited system suffered from five interrelated architectural deficiencies: a blocking checkerboard detection call that throttled the vision pipeline from 10~Hz to 3~Hz; a systematic area-based LED anchor heuristic that always selected ceiling reflections over the true beacon, yielding a 0\% detection rate; a checkerboard landmark detector effectively disabled by a 1-second gating timer; an instant-trip heartbeat failsafe caused by stale state persisting across browser sessions; and non-deterministic network interface selection that silently discarded all operator commands. Individually resolvable, their combination rendered the autonomous state machine entirely untestable.

The work resolved these deficiencies through systematic re-engineering across seven dimensions: a multiprocess vision architecture isolating OpenCV workloads in a \code{spawn}-context subprocess with bounded IPC queues; a density-based LED clustering algorithm with a formally provable correctness guarantee under an empirically validated isolation hypothesis; promotion of \code{findChessboardCornersSB} to the primary detection path in all code paths; a session-aware heartbeat stale guard; deterministic \code{ROS\_IP} configuration; a complete redesign of the autonomous finite-state machine into six states with dual-frame (local/world) odometry and persistent map markers; and a browser-based operator cockpit built around a NASA/FIT-inspired glassmorphic design language, a \code{requestAnimationFrame}-decoupled Canvas~2D vision overlay, an $O(1)$ ring-buffer trajectory map, and a six-layer safety architecture.

Validation demonstrates a 98.7\% LED beacon detection rate and 100\% checkerboard detection rate (both with 95\% confidence intervals exceeding the 95\% target objective at $p<0.05$), a reduction in vision processing latency from 200--330~ms to 5--20~ms, zero false failsafe trips across repeated reconnection cycles, and a fully autonomous 45-minute, 87.3~m patrol mission with five beacons logged, zero operator interventions, and a 14~cm world-frame return-to-base position error. A reliability analysis combining Markov chain modeling and fault tree analysis estimates the probability of an uncommanded, unmonitored motion event at $\sim4\times10^{-17}$ per operational hour, nine orders of magnitude below the IEC~61508 SIL~4 threshold, though this estimate is offered as a qualitative robustness indicator rather than a formal certification claim.

The resulting system is deployed and operational in the FIT robotics laboratory and serves as the reference implementation for the LEO Rover platform within the broader Carolus comparative study of the LIMO Rover, LEO Rover, and DJI RoboMaster~S1 platforms.
